\section{Introduction}
\label{sec:introduction}

Authenticated encryption with associated data (AEAD) is often treated in
practice as if its tag were a compact commitment to the full encryption context:
a short string that binds the key, nonce, associated data, and message. One
source of this intuition is that many practitioners first meet message
authentication through hash-based MACs such as HMAC, whose tags look and behave
much like keyed digests. Standard modern AEADs do not provide this full
commitment property at all. AES-GCM and ChaCha20-Poly1305, for example, are not
key-committing: an adversary can craft a single ciphertext that decrypts to
valid plaintexts under two or more keys \cite{DGRW18, ADGKLS22}. The cause is
structural. Their authenticity rests on algebraic MACs whose outputs can often
be made to validate under a second context, whereas a digest is expected to
resist both collisions and preimages.

The resulting attacks are not edge cases in the formalism; they are failures of
the compact-commitment intuition that users and protocols naturally assign to
tags. The ``invisible salamander'' attack defeats Facebook's attachment
franking scheme \cite{DGRW18}; related constructions produce ciphertexts that
open to different well-formed files under different keys \cite{ADGKLS22};
partitioning oracles recover passwords from systems built on non-committing AEAD
\cite{LGR21}; and context discovery attacks break \textsf{GCM}, \textsf{CCM},
\textsf{EAX}, \textsf{SIV}, and \textsf{OCB3} by exploiting authentication
layers that are not preimage-resistant \cite{MLGR23}. These attacks differ in
surface form, but all expose the same mismatch: the tag is used as if it were a
full-context commitment, while the mode only proves ordinary authenticity.

The AEAD literature has only recently begun to model this expectation directly.
Compactly committing AEAD was introduced for message franking \cite{GLR17}, and
the encryptment approach of \cite{DGRW18} builds commitment into a dedicated
primitive. The committing hierarchy of \cite{BH22} makes the collision-style
requirement explicit: no two distinct contexts should open the same ciphertext
and tag. Context discoverability \cite{MLGR23} captures the preimage-facing
counterpart: given a target ciphertext, it should be infeasible to find a
decrypting context. These notions explain when a ciphertext can have another
valid opening, or any valid opening at all, but they are still phrased as AEAD
games over ciphertexts. The practitioner's assumption is more tag-centric: the
tag itself is treated as a compact commitment string. We therefore analyze the
tag with hash notions directly, and then recover the AEAD commitment notions as
game-specific consequences of the same tag analysis.

\TW{} (\emph{TreeWrap}, with $128$ denoting the $128$-bit security target)
takes the digest intuition literally. If an AEAD tag is meant to behave like a
compact digest, the mode should be shaped like a modern high-performance digest
rather than like a serial stream cipher with an algebraic MAC. The relevant
performance lesson from modern hashing is that large inputs are processed by
splitting them into chunks, evaluating independent work in parallel, and
aggregating short chaining values into a final root value.
KangarooTwelve \cite{BDPVVV16} is the SHA-3-family exemplar of this pattern: it
hashes fixed-size chunks independently under \Keccakp{}, absorbs their chaining
values into a root node, runs leaves across SIMD lanes and cores, and keeps
memory bounded independently of the input length.

\TW{} asks what this tree looks like when the sponge nodes are replaced by
keyed duplexes that encrypt while they absorb. A serial root duplex, initialized
with the key $K$ and nonce $N$, absorbs the associated data, encrypts the first
message chunk, absorbs a tag from each remaining chunk, and emits the final root
tag. Each remaining chunk is processed by an independent leaf duplex keyed by
$(K, N, j)$, where $j$ is the chunk index; the leaf encrypts its chunk and
returns a leaf tag to the root. Thus every input reaches the root directly or
through a leaf tag, while every duplex state remains bounded at $1600$ bits. As
in KangarooTwelve's \emph{kangaroo hopping}, a single-chunk message uses only
the root duplex and pays no tree overhead (\Cref{sec:latency}).

This gives a mode that is not merely parallelizable but scalable. The leaf
decomposition and transcript are canonical, so there is only one ciphertext and
tag for a given input, but the number of leaves evaluated at once is left to the
implementation. A scalar implementation, an AVX2 implementation, an AVX-512
implementation, a NEON implementation, and a multicore implementation all
produce identical ciphertexts and tags. The degree of parallelism is therefore
not a mode parameter and is not visible on the wire. This separates \TW{} from
prior duplex-based AEADs with parallel variants, where the parallelism degree is
fixed by the mode instance or otherwise shared by the communicating parties.

Because \TW{} adapts the KangarooTwelve tree to encryption, the tag emitted by
the root is naturally the output of a keyed tree hash. We formalize this by
considering the wrapper $\HTW$ that returns the root tag, and prove that $\HTW$
is a secure hash in the sense of \cite{RS04}: collision-, second-preimage-, and
preimage-resistant as a function of the full $(K, N, A, M)$ input. The $256$-bit
root tag is therefore a compact digest-like commitment to the full context. The
familiar AEAD commitment notions are recovered from the same root tag candidate
analysis: full-input $s$-way \textsf{CMTDs-4} \cite{BH22} is the
collision-facing AEAD consequence, while \cite{MLGR23} context discoverability
$\mathsf{CDY}^\ast$ is the preimage-facing consequence, with an additional
hidden-component guessing term. The proofs use the structure of the AEAD games
to keep the stated bounds tight.

We also prove the standard AEAD security of the mode. In the public permutation
model, \TW{} has a concrete multi-user indistinguishability bound in the sense
of \cite{BT16}, with all-in-one authenticated encryption in the sense of
\cite{RS06} as its single-user corollary. Each bound is proved first against an
intermediate random oracle and then transferred to the real permutation by the
flat sponge lift of \cite{BDPV08} (\Cref{app:flat-sponge-lift}). At the
implemented parameters ($k = c = \tauroot = \tauleaf = 256$), the bounds meet a
$128$-bit security target for adversaries running up to $\Ntot \leq 2^{63}$
\Keccakp{} calls---on the order of $2^{71}$ bytes of processed data.

The mode's design-level parallelism translates into measured implementation
performance. Optimized \texttt{amd64} (AVX-512 and AVX2) and \texttt{arm64}
(NEON with the SHA-3 extension) implementations evaluate leaf duplexes in
batches while preserving the canonical transcript. On contemporary hardware,
\TW{} exceeds $10$\,Gbit/s and reaches $50$--$75\%$ of hardware-accelerated
AES-128-GCM
(\Cref{sec:performance}), while using one permutation for confidentiality,
authentication, and commitment.

\subsection{Our Contributions}
\label{sec:contributions}

\begin{itemize}
\item \textbf{A scalable tree AEAD mode.} We specify \TW{}, a nonce-based AEAD
  with associated data built as a two-level tree of strict \cite{BDPV11}
  duplexes over \Keccakp{} (\Cref{sec:tw128}). The mode fixes the root/leaf
  transcript but delegates the number of leaves evaluated simultaneously to the
  implementation, so scalar and vectorized implementations interoperate without
  changing ciphertexts or tags.

\item \textbf{Digest-like commitment.} We prove that the wrapper $\HTW$
  returning the root tag is a secure hash of the full $(K, N, A, M)$ input in
  the sense of \cite{RS04}: collision-, second-preimage-, and
  preimage-resistant (\Cref{thm:coll,thm:asec,thm:epre}). Thus the
  authentication tag behaves like the compact digest practitioners often expect
  an AEAD tag to be.

\item \textbf{AEAD commitment consequences.} We instantiate the same root tag
  analysis in the existing AEAD vocabulary. The all-bodies-equal structure of
  \textsf{CMTDs-4} gives the full-input $s$-way committing bound without the
  leaf-collision slack of the general hash collision theorem \cite{BH22}
  (\Cref{cor:cmtds4}), and the remaining \textsf{CMT}/\textsf{CMTD} notions
  for $\ell \in \{1, 3, 4\}$ follow from the committing hierarchy and direct
  source-game arguments (\Cref{cor:committing-combined}). The same fixed-target
  root-preimage analysis, plus hidden-component guessing, gives
  $\mathsf{CDY}^\ast$ context discoverability security \cite{MLGR23}
  (\Cref{thm:cdy-full}).

\item \textbf{Multi-user and AE security.} We give a concrete multi-user
  indistinguishability bound \cite{BT16} for $D$ users (\Cref{thm:mu-ind}).
  The $D = 1$ case yields a concrete all-in-one AE bound \cite{RS06} in the
  public permutation model (\Cref{cor:ae}). The proof proceeds through a
  dedicated \TW{} random oracle model and is transferred to the public
  permutation setting by the \cite{BDPV08} flat sponge lift.

\item \textbf{Optimized implementations.} We describe and evaluate vectorized
  \texttt{amd64} (AVX-512 and AVX2) and \texttt{arm64} (NEON with the SHA-3
  extension) implementations whose batched leaf cores realize the mode's
  scalable parallelism in practice (\Cref{sec:performance}).
\end{itemize}

\subsection{Related Work}
\label{sec:related}

\paragraph{Committing AEAD.}
The study of committing authenticated encryption was initiated by \cite{GLR17},
who introduced compactly committing AEAD for message franking, and by
\cite{DGRW18}, who built it from a dedicated primitive (encryptment).
\cite{BH22} organized the collision-style notions into the
\textsf{CMT-1}/\textsf{3}/\textsf{4} and \textsf{CMTD} hierarchy and gave the
\textsf{UtC}, \textsf{RtC}, and \textsf{HtE} transforms; \cite{CR22} gave the
\textsf{CTX} transform, which commits a nonce-based AE scheme to its full
context by adding a hash. These transforms are efficient and broadly
applicable, but they treat commitment as an external layer. \TW{} instead makes
the authentication tag itself the hash-like commitment: the root tag is the
compact commitment, and the root tag analysis gives both the \cite{BH22}
full-input committing consequence and the \cite{MLGR23} context discoverability
bound. Closest in spirit is the recent work
of \cite{DHMVV24}, which also builds nonce-based AE natively on SHA-3 functions
and obtains \textsf{CMT-4} from collision resistance; it targets session-based
communication through a sequential duplex chain, whereas \TW{} is a single
scalable tree for standalone AEAD
(\Cref{sec:compare-commit}).

\paragraph{Sponge and duplex AEADs.}
Permutation-based authenticated encryption built on the sponge and duplex
constructions of \cite{BDPV08, BDPV11} is well established. Keyak
\cite{BDPVV16} uses the Motorist mode over round-reduced \Keccak{} and already
admits parallel duplex instances; Ascon \cite{DEMS21} and Xoodyak
\cite{DHPVV20} are serial duplex AEADs over smaller permutations. \TW{} shares
the single-permutation, constant-state duplex foundation but differs in where
parallelism lives. Keyak fixes the degree of parallelism as a mode parameter,
while Ascon and Xoodyak expose no mode-level parallelism. In \TW{}, the
decomposition into indexed leaves is fixed by the transcript, but the number of
leaves evaluated at once is chosen by the implementation and is invisible to the
wire format (\Cref{sec:compare-sponge}).

\paragraph{Tree hashing.}
\TW{}'s two-level shape is inspired by KangarooTwelve \cite{BDPVVV16}, a tree
hash over the same permutation. \TW{} adopts its choice of permutation and
final-node/leaf split and adapts them to authenticated encryption: each leaf
chunk is encrypted and tagged independently, leaf tags replace unkeyed chaining
values, and the root is keyed by $(K, N)$. The construction also inherits
KangarooTwelve's short-message behavior: inputs fitting in one chunk use only
the root duplex. Because \TW{} has a fixed two-level topology with independently
keyed, indexed leaves and phase-separated duplex calls, it does not need
KangarooTwelve's Sakura tree encoding \cite{BDPV14} to make its transcripts
unambiguous (\Cref{sec:tw128}).

\subsection{Paper Organization}
\label{sec:organization}

\Cref{sec:prelims} recalls the notation, sponge and duplex constructions,
nonce-based AEAD security, committing security \cite{BH22}, context
discoverability \cite{MLGR23}, and the hash notions of \cite{RS04}.
\Cref{sec:tw128} specifies the \TW{} construction and parameters.
\Cref{sec:security-model} fixes the public permutation games, intermediate
random oracle model, resource accounting, and loss terms.
\Cref{sec:aead-security} proves multi-user indistinguishability and all-in-one
AE security. \Cref{sec:root-tag-hash-security} proves root-tag hash security,
committing security, and context discoverability, and
\Cref{sec:concrete-security} summarizes the concrete bounds.
\Cref{sec:performance} reports implementation performance, and
\Cref{sec:discussion} compares \TW{} to prior AEAD and committing constructions
and discusses limitations. The appendices collect structural lemmas, the flat
sponge lift, vectorized kernel details, and test vectors.
